\section*{Capítulo \ref{ch:g3:numbers} --- Números hasta 10\,000}

\begin{solution}{exo:g3:numbers:1}
$2\,345$; \quad $6\,050$; \quad $9\,001$.
\end{solution}

\begin{solution}{exo:g3:numbers:2}
$1\,208$: mil doscientos ocho. \quad
$4\,070$: cuatro mil setenta. \quad
$9\,999$: nueve mil novecientos noventa y nueve.
\end{solution}

\begin{solution}{exo:g3:numbers:3}
Dígito de las decenas: $8$. Dígito de miles: $7$. Cientos contenidos en total:
$73$ (porque $7\,382 = 73$ cientos $+ 82$, es decir.\
$73 \times 100 + 82$).
\end{solution}

\begin{solution}{exo:g3:numbers:4}
$5\,631 = 5$ miles $+ 6$ cientos $+ 3$ decenas $+ 1$ unidad.

$2\,046 = 2$ miles $+ 0$ cientos $+ 4$ decenas $+ 6$ unidades.

$8\,500 = 8$ miles $+ 5$ cientos.
\end{solution}

\begin{solution}{exo:g3:numbers:5}
$456 < 465$ (decenas: $5 < 6$); \quad
$1\,203 > 989$ (cuatro dígitos ganan a tres); \quad
$6\,090 > 6\,009$ (decenas: $9 > 0$); \quad
$9\,999 < 10\,000$.
\end{solution}

\begin{solution}{exo:g3:numbers:6}
$353 < 503 < 530 < 3\,500 < 5\,033$.
\end{solution}

\begin{solution}{exo:g3:numbers:7}
$150$ se sienta a medio camino entre $100$ y $200$; $480$ justo antes
$500$; $820$ justo después $800$; $990$ casi en $1\,000$.
\end{solution}

\begin{solution}{exo:g3:numbers:8}
\omterm{def:g3:numbers:evenodd}{par}: $34$, $70$, $6\,548$ (último dígito $4$, $0$, $8$). \omterm{def:g3:numbers:evenodd}{impar}: $57$,
$2\,481$ (último dígito $7$, $1$).
\end{solution}

\begin{solution}{exo:g3:numbers:9}
De $2\,970$: $2\,980$, $2\,990$, $3\,000$, $3\,010$, $3\,020$.

De $860$: $960$, $1\,060$, $1\,160$, $1\,260$.

De $9\,996$: $9\,997$, $9\,998$, $9\,999$, $10\,000$ --- el último
El paso cruza a un número de cinco dígitos: ¡el límite de todo nuestro capítulo!
\end{solution}

\begin{solution}{exo:g3:numbers:10}
Más grande: $7\,530$ (dígitos en orden decreciente). Pequeñísimo: $3\,057$
(primero el dígito más pequeño distinto de \omterm{def:g1:counting:zero}{cero}, luego el resto aumenta --- el $0$
justo después del $3$).
\end{solution}

\begin{solution}{exo:g3:numbers:11}
Más grande que $800$: el dígito de las centenas es $8$. \omterm{def:g3:numbers:evenodd}{par}: el dígito de las unidades
debe ser $2$ (usando cada uno de $2$, $5$, $8$ una vez). Entonces el número es
$852$.
\end{solution}
